%表紙
\newpage
\chapter{教務主事}
\begin{enumerate}[label=第\arabic*条　]
  \item この章は、憲章の趣旨に基づき、教務主事がその職務を行うに当たって拠るべき基本理念を定めるものとする。
    \begin{enumerate}[label=\arabic*　,start=2,leftmargin=*]
      \item 教務主事の職務に関する具体的な手続については、別に定める。
    \end{enumerate}
  \item 教務主事は、学園の管理運営に関する事務を担い、学生が憲章の定める探究活動を継続的に行うことのできる環境を維持する責務を負う。
  \item 教務主事は、その職務が、学生の探究活動を可能にし、これを支えるためのものであることを常に自覚しなければならない。
    \begin{enumerate}[label=\arabic*　,start=2,leftmargin=*]
      \item 教務主事は、活動の形式又は対象を問わず、学生の探究活動を等しく支援するものとし、その内容の優劣を判断する立場にないことを認識しなければならない。
    \end{enumerate}
  \item 教務主事の地位は、職務の遂行のために付与された機能であって、学生に対する人格的な優越を意味するものではない。
    \begin{enumerate}[label=\arabic*　,start=2,leftmargin=*]
      \item 教務主事は、学生としての資格を保持したまま職務に当たる者であり、他の学生と対等な立場において交わるものとする。
    \end{enumerate}
  \item 教務主事が学生の活動に対して行う介入は、憲章の定める探究の自由を不当に制限することのないよう、必要な最小限度にとどめなければならない。
    \begin{enumerate}[label=\arabic*　,start=2,leftmargin=*]
      \item 教務主事は、複数の手段を採り得る場合には、より制限の程度が小さく、かつ回復の容易な手段を優先して選択するものとする。
      \item 教務主事は、その権限を、学園規則に定める事由がある場合に限り行使するものとし、これを他の目的のために用いてはならない。
    \end{enumerate}
  \item 教務主事は、その職務上の判断について、理由を述べることができる状態を保たなければならない。
    \begin{enumerate}[label=\arabic*　,start=2,leftmargin=*]
      \item 教務主事は、学生の権利利益に影響を及ぼす判断を行ったときは、当該学生に対し、その理由を明らかにするものとする。
    \end{enumerate}
  \item 教務主事は、その職務が構成員の交代を超えて継承されるべきものであることを自覚し、判断の基準及び経緯を記録に残すよう努めるものとする。
    \begin{enumerate}[label=\arabic*　,start=2,leftmargin=*]
      \item 教務主事は、自己の担当する事務が特定の個人に依存することのないよう、手続の明文化に努めるものとする。
    \end{enumerate}
  \item 教務主事が複数置かれるときは、相互に情報を共有し、協力して職務に当たるものとする。
  \item この章の実施に関し必要な事項は、学園長が別に定める。
\end{enumerate}
